\begin{abstract}
This companion records a selection of substantial results in the Lean 4 formalization accompanying \emph{Models for dagger $(\infty,1)$-categories I: Dagger simplicial sets and unitary cores}.
It includes the four main theorems of Part I and supporting results in dagger category theory, model category theory, simplicial homotopy theory, and equivariant constructions.
Each entry gives a mathematical statement, the corresponding Lean declarations, source locations, and a brief account of the construction or reused infrastructure.
The formalization was developed with AI assistance.
\end{abstract}

\paragraph{AI-generated text.}
The prose in this catalogue, including the English translation and summaries of the formalized results, was generated by AI agents from the manuscript and Lean sources.
It should be read alongside the cited declarations and the paper.

\section*{Scope and conventions}
\addcontentsline{toc}{section}{Scope and conventions}
The entries are selected by mathematical content rather than by the Lean keywords \texttt{lemma}, \texttt{theorem}, \texttt{def}, or \texttt{instance}.
Constructions of model structures, adjunctions, and equivalences are included as existence results.
Routine calculations, individual lifting steps, type conversions, and closure checks are grouped with the larger results they support.
This is a catalogue, not a replacement for the proof source or an exhaustive list of declarations.

\begin{center}
\begin{tabular}{lr}
\toprule
Area & Entries\\
\midrule
Main theorems A--D of Part I & 4\\
Dagger categories and simplicial constructions & 19\\
Model category theory and concrete model structures & 10\\
Simplicial homotopy theory and group actions & 7\\
\midrule
Total & 40\\
\bottomrule
\end{tabular}
\end{center}
These forty entries include related constructions and alternative descriptions of the same structures.
They should not be read as forty independent new mathematical theorems.

\paragraph{Notation from Part I.}
We retain the notation of Part I for categories, functors, generating maps, and model structures.
The categories of simplicial sets, simplicial categories, dagger simplicial sets, dagger simplicial categories, and dagger simplicial graphs are denoted by
$\sSet$, $\sCat$, $\sSet^{\dag}$, $\sCat^{\dag}$, and $\sGph^{\dag}$, respectively.
The category $\sSet^{\AInv}$ consists of anti-involutive simplicial sets; its involution need not fix vertices.
The relevant adjunctions include
\begin{align}
\FdagSSet &: \sSet\rightleftarrows\sSet^{\dag}:\UdagSSet,\\
\FCatDag &: \sGph^{\dag}\rightleftarrows\sCat^{\dag}:\UCatDag,\\
\frakC_{\dag} &: \sSet^{\dag}_{\Joyal}\rightleftarrows\sCat^{\dag}_{\Bergner}:N_{\dag}.
\end{align}
Here $\FdagSSet$ freely adds a dagger to a simplicial set, whereas $\FCatDag$ is the free-category functor on dagger simplicial graphs.
The functor forgetting a dagger on a simplicial category is $\UdagSCat$.
We use $I_{\dag},J_{\dag},I^{\sSet}_{\dag}$ for the generating families and $W_{\dag},W^{\sSet}_{\dag}$ for the two dagger weak-equivalence classes.
The unitary core is $\calU(X)$, the groupoid core is $X^{\simeq}$, and the zero-skeleton is $X^{(0)}=\sk_0X$.
The reversal simplex category is $\prism_{\rev}$.
We write $C_2=\{1,\omega\}$ with $\omega^2=1$, and distinguish the simplicial set $E_{\bullet}C_2=N(\widetilde{C_2})$ from its topological realization $EC_2=|E_{\bullet}C_2|$.

For a family of maps $I$, the notation $I\text{-}\inj$ denotes its right lifting class, $I\text{-}\cell$ the relative cell complexes, and $I\text{-}\cof$ their retract closure.
We also use ${}^{\square}I$ and $I^{\square}$ for general left and right orthogonality classes.
Unless another model structure is specified, weak equivalences and fibrations of simplicial sets refer to the Kan--Quillen model structure, as in Part I.
Finitely presentable objects are defined by preservation of filtered colimits by their Hom functors.

Universe conditions are stated in the text rather than appended to every category symbol.
The notation $\Set_u$ means the category of sets in the value universe corresponding to Lean's \texttt{Type u}.
Auxiliary constructions not named in Part I are given explicit supplementary notation.
For example, the right-cone cosimplicial object used in the implementation is denoted by $Q_{\mathrm{rc}}$; it is distinct from the self-conjugate simplex $Q^n$ of Part I.
Lean declaration names are retained verbatim in the correspondence fields.

\paragraph{Provenance.}
A project-specific construction or proof may use substantial Mathlib infrastructure.
Applications of general presheaf or presentability theorems are identified as such.
Results reused from Joël Riou's \texttt{TopCatModelCategory} are credited separately and are not counted as new proofs by the author.
The selection describes formalization work, without making a claim of mathematical novelty for each entry.

\paragraph{Verification.}
The entries refer to the saved source validated on October 18, 2026.
Their mathematical wording has been checked against the corresponding Lean types.
Kernel checking of a Lean proof and correspondence between its statement and the intended mathematics are separate questions.
The final section records the validation evidence and its scope.
